% Support de soutenance Ricochet — Beamer.
%
% Structure imposée par le brief, et budget de temps respecté par le nombre de
% diapositives :
%
%   | Partie                                | Temps  | Diapositives |
%   | Approches de modélisation testées     | 10 min |      8       |
%   | Fonctionnalités du système dans l'app |  6 min |      5       |
%   | Architecture technique cible          |  2 min |      2       |
%   | Démonstration de l'application        |  2 min |      1       |
%
% Compilation (deux passes pour la table des matières) :
%   latexmk -pdf -xelatex docs/presentation.tex
% ou :
%   xelatex -output-directory=docs docs/presentation.tex   (deux fois)
%
% XeLaTeX plutôt que pdfLaTeX : le texte contient des caractères (★, ×, —) et une
% typographie française qui passent mal en encodage historique.

\documentclass[aspectratio=169, 11pt]{beamer}

\usepackage{fontspec}
\usepackage{polyglossia}
\setdefaultlanguage{french}
\usepackage{booktabs}
\usepackage{array}
\usepackage{tikz}
\usetikzlibrary{arrows.meta, positioning}

% --- Thème sobre : pas de fioritures, le contenu porte le propos -------------
\usetheme{default}
\usecolortheme{default}
\setbeamertemplate{navigation symbols}{}
\setbeamertemplate{itemize items}[circle]

\definecolor{primary}{RGB}{31,58,95}
\definecolor{accent}{RGB}{46,134,193}
\definecolor{accent2}{RGB}{230,126,34}
\definecolor{vert}{RGB}{39,174,96}
\definecolor{gris}{RGB}{127,140,141}
\definecolor{clair}{RGB}{244,246,247}

\setbeamercolor{frametitle}{bg=primary, fg=white}
\setbeamercolor{title}{fg=white}
\setbeamercolor{structure}{fg=accent}
\setbeamercolor{alerted text}{fg=accent2}
\setbeamerfont{frametitle}{size=\large, series=\bfseries}
\setbeamerfont{framesubtitle}{size=\small}
\setbeamercolor{framesubtitle}{fg=gris}

\setbeamertemplate{footline}{%
  \hfill\usebeamerfont{footline}\color{gris}\scriptsize
  \insertframenumber\,/\,\inserttotalframenumber\hspace*{2ex}\vspace*{1ex}}

\newcommand{\pointfort}[1]{\textcolor{accent2}{\textbf{#1}}}

% Bloc « à retenir » en bas de diapositive
\newcommand{\aretenir}[1]{%
  \vfill
  \begin{beamercolorbox}[wd=\textwidth, sep=2ex, rounded=true]{}%
    \color{primary}\small #1
  \end{beamercolorbox}}

\title{\Huge Ricochet}
\subtitle{Système de recommandation d'articles — MVP}
\date{}

\begin{document}

% ═══════════════════════════════════════════════════════════ Ouverture ══════
{\setbeamercolor{background canvas}{bg=primary}
\begin{frame}[plain]
  \vfill
  {\color{white}\Huge\bfseries Ricochet}\\[0.6em]
  {\color{white!75!primary}\Large Système de recommandation d'articles — MVP}\\[1.5em]
  \textcolor{accent2}{\rule{0.25\textwidth}{2pt}}
  \vfill
  {\color{white!70!primary}\small Approches testées \textperiodcentered\ Fonctionnalités
   \textperiodcentered\ Architecture cible \textperiodcentered\ Démonstration}
\end{frame}}

\begin{frame}{Contexte}
  \framesubtitle{Ce qui est demandé, avec quelles données}
  \begin{itemize}
    \item Proposer \pointfort{5 articles} à chaque lecteur, pour l'inciter à lire
          davantage.
    \item Pas de données propriétaires : jeu public Globo.com —
          \textbf{2 988 181} clics, \textbf{322 897} lecteurs,
          \textbf{364 047} articles.
    \item[] \vspace{0.3em}
    \item \pointfort{Contrainte majeure : le jeu est anonymisé.}
      \begin{itemize}
        \item aucun titre, aucun texte : un article n'est qu'un vecteur de
              250 dimensions et un identifiant ;
        \item aucune note, aucun avis : seulement des clics, donc un signal
              implicite ;
        \item conséquence — tout se fait sans lire le contenu, et l'interface ne
              peut afficher que des identifiants.
      \end{itemize}
  \end{itemize}
\end{frame}

% ═══════════════════════ PARTIE 1 — Modélisation (10 min, 8 diapos) ════════
\begin{frame}{Partie 1 — Les cinq approches testées}
  \framesubtitle{10 minutes \textperiodcentered\ comment chacune classe les articles}
  \centering
  \small
  \begin{tabular}{@{}lll@{}}
    \toprule
    \textbf{Approche} & \textbf{Principe} & \textbf{Bibliothèque} \\
    \midrule
    Popularité       & articles les plus lus sur une fenêtre     & --- \\
    Contenu          & profil = moyenne des embeddings, cosinus  & scikit-learn (ACP) \\
    Collaboratif ALS & factorisation lecteur $\times$ article    & implicit \\
    Collaboratif SVD & prédiction de notes dérivées des clics    & Surprise \\
    Mixte            & combinaison sur les 5 places              & --- \\
    \bottomrule
  \end{tabular}
  \aretenir{La popularité sert de \pointfort{référence à battre} : elle ne
    personnalise rien, donc toute approche plus complexe doit faire mieux qu'elle.}
\end{frame}

\begin{frame}{Approche 1 — Filtrage par le contenu}
  \framesubtitle{Similarité des embeddings d'articles}
  \centering
  \begin{tikzpicture}[node distance=8mm, every node/.style={font=\scriptsize},
      bloc/.style={draw=primary, thick, rounded corners, text width=2.6cm,
                   align=center, minimum height=1cm, text=white}]
    \node[bloc, fill=accent]  (a) {Articles lus\\par le lecteur};
    \node[bloc, fill=accent2, right=of a] (b) {Profil = moyenne\\des vecteurs};
    \node[bloc, fill=primary, right=of b] (c) {Cosinus vers\\le catalogue};
    \draw[-{Latex[length=2mm]}, gris, thick] (a) -- (b);
    \draw[-{Latex[length=2mm]}, gris, thick] (b) -- (c);
  \end{tikzpicture}

  \vspace{1em}
  \scriptsize
  \begin{columns}[t]
    \begin{column}{0.48\textwidth}
      \textcolor{vert}{\textbf{Avantages}}
      \begin{itemize}\itemsep2pt
        \item fonctionne dès le 1\textsuperscript{er} clic, sans ré-entraînement ;
        \item gère un article nouveau immédiatement (il a un vecteur) ;
        \item couverture du catalogue la plus large : \textbf{0,27 \%} contre 0,003 \% ;
        \item personnalisation \textbf{96 \%}.
      \end{itemize}
    \end{column}
    \begin{column}{0.48\textwidth}
      \textcolor{red!70!black}{\textbf{Inconvénients}}
      \begin{itemize}\itemsep2pt
        \item précision faible : HitRate@5 = \textbf{0,0200} ;
        \item n'exploite pas l'intelligence collective ;
        \item sans restriction aux articles récents, propose des articles vieux de
              plusieurs années que personne ne lit.
      \end{itemize}
    \end{column}
  \end{columns}
\end{frame}

\begin{frame}{Approche 2 — Collaboratif ALS (\texttt{implicit})}
  \framesubtitle{Factorisation pour feedback implicite}
  \small
  \begin{itemize}
    \item Traite toutes les cases vides comme des \textbf{négatifs pondérés} :
          formulation adaptée à des clics, pas à des notes.
    \item Réglages mesurés sur validation — les deux pointent dans le même sens :
      \begin{itemize}
        \item fenêtre d'entraînement \textbf{72 h} (0,0200) plutôt que tout
              l'historique (0,0120) ;
        \item \textbf{16 facteurs} latents (0,0345) plutôt que 100 (0,0125).
      \end{itemize}
    \item \pointfort{Moins de données, plus récentes, et un modèle plus petit}
          valent mieux qu'un gros modèle appris sur dix jours d'actualité.
    \item Résultat sur test : HitRate@5 = \textbf{0,0415}, personnalisation 95 \%.
  \end{itemize}
  \aretenir{Limite structurelle : un article publié il y a deux heures n'a presque
    aucune co-lecture — le collaboratif ne peut pas le proposer.}
\end{frame}

\begin{frame}{Approche 3 — Collaboratif SVD (\texttt{Surprise})}
  \framesubtitle{Trois définitions de note, trois résultats}
  \centering
  \scriptsize
  \begin{tabular}{@{}lccl@{}}
    \toprule
    \textbf{Définition de la note} & \textbf{HitRate@5} & \textbf{Person.} &
    \textbf{Verdict} \\
    \midrule
    Étoiles de l'article (log des clics) & 0,0000 & 12 \% &
      reproduit la popularité \\
    Binaire + 1 négatif échantillonné   & 0,0080 & 46 \% & insuffisant \\
    Binaire + 4 négatifs échantillonnés & 0,0755\,$^*$ & 34 \% &
      meilleure variante \\
    \bottomrule
  \end{tabular}

  \vspace{0.8em}
  \raggedright\scriptsize
  \begin{itemize}\itemsep3pt
    \item Même bibliothèque, même algorithme : c'est la \pointfort{formulation du
          problème} qui change tout, pas l'outil.
    \item Les étoiles échouent parce que la note ne dépend que de l'article : le
          biais article capte tout le signal, et le modèle apprend « quels articles
          sont lus », pas « qui lit quoi ».
    \item $^*$ mesuré sur validation ; ce réglage \textbf{ne s'est pas reproduit}
          sur le test (0,0160) — d'où l'utilité d'une période de test intacte.
  \end{itemize}
\end{frame}

\begin{frame}{Comment nous avons su que les premiers chiffres étaient faux}
  \framesubtitle{Découpage temporel 60 / 20 / 20}
  \small
  \begin{itemize}
    \item Première évaluation : modèles entraînés sur 100 \% des données, puis
          évalués sur ces mêmes données. L'ALS affichait HitRate@5 = 0,2300.
    \item Avec un découpage temporel correct, il tombe à \textbf{0,0250} —
          \pointfort{un facteur 42}, invisible sans découpage.
    \item Temporel et non aléatoire : sur un flux d'actualité, un découpage
          aléatoire laisse le modèle apprendre des clics \emph{postérieurs} à ceux
          qu'il doit prédire.
    \item Le réglage se fait sur la validation ; le test ne sert qu'à la mesure
          finale, une seule fois.
  \end{itemize}

  \vspace{0.5em}
  \centering
  \begin{tikzpicture}[every node/.style={font=\scriptsize, text=white,
        minimum height=0.9cm, align=center}]
    \node[fill=accent,  minimum width=6cm]   (a) {60 \% entraînement};
    \node[fill=accent2, minimum width=2.6cm, right=0pt of a] (b) {20 \% validation};
    \node[fill=primary, minimum width=2.6cm, right=0pt of b] (c) {20 \% test};
  \end{tikzpicture}

  {\scriptsize\color{gris} le temps s'écoule de gauche à droite — aucune fuite du
   futur vers le passé}
\end{frame}

\begin{frame}{Résultat principal : la fraîcheur pèse plus que le modèle}
  \framesubtitle{Même algorithme, deux fenêtres de comptage}
  \centering
  \small
  \begin{tabular}{@{}lrr@{}}
    \toprule
    \textbf{Fenêtre de comptage des clics} & \textbf{Candidats} &
    \textbf{HitRate@5} \\
    \midrule
    1 heure                      &    791 & \textbf{0,2190} \\
    6 heures                     &  2 657 & 0,1480 \\
    24 heures                    &  6 777 & 0,1060 \\
    Tout l'historique (240 h)    & 29 118 & \textbf{0,0075} \\
    \bottomrule
  \end{tabular}

  \vspace{0.8em}
  \raggedright\small
  \begin{itemize}\itemsep3pt
    \item \pointfort{Facteur 250} entre la première et la dernière ligne, sans
          changer une ligne d'algorithme.
    \item Explication : \textbf{48 \%} des articles lus pendant la période évaluée
          n'existaient pas encore pendant l'entraînement.
    \item Conséquence produit : recalculer la fenêtre en continu compte davantage
          que raffiner le modèle.
  \end{itemize}
\end{frame}

\begin{frame}{Quatre critères, pas un seul}
  \framesubtitle{La précision seule désigne toujours la popularité}
  \centering
  \scriptsize
  \begin{tabular}{@{}llrr@{}}
    \toprule
    \textbf{Critère} & \textbf{Ce qu'il mesure} & \textbf{Popularité} &
    \textbf{Contenu} \\
    \midrule
    HitRate@5        & un article recommandé a été lu      & 0,2525 & 0,0200 \\
    Couverture       & part du catalogue exposée           & 0,003 \% & 0,27 \% \\
    Personnalisation & différence entre deux listes        & 4 \%   & 96 \% \\
    Recall@5         & part du futur devinée (plafond 0,76)& 0,0555 & 0,0034 \\
    \bottomrule
  \end{tabular}

  \vspace{0.8em}
  \raggedright\small
  \begin{itemize}\itemsep3pt
    \item La popularité est 12 fois plus précise, mais expose \textbf{11 articles
          sur 364 047} : un éditeur y perd son catalogue.
    \item C'est le quatrième critère qui a \pointfort{disqualifié le SVD sur
          étoiles} : personnalisation 12 \%, soit \emph{moins} que la popularité —
          il la reproduisait en plus grossier.
  \end{itemize}
\end{frame}

\begin{frame}{Résultat final sur la période de test}
  \framesubtitle{Chaque approche dans sa meilleure configuration}
  \centering
  \scriptsize
  \begin{tabular}{@{}lrrr@{}}
    \toprule
    \textbf{Configuration} & \textbf{HitRate@5} & \textbf{Couverture} &
    \textbf{Person.} \\
    \midrule
    Popularité 1 h                        & \textbf{0,2525} & 0,003 \% &  4 \% \\
    \textbf{Mixte : 4 populaires + 1 contenu} & \textbf{0,2500} & 0,116 \% & 22 \% \\
    ALS (72 h, 16 facteurs)               & 0,0415 & 0,015 \% & 95 \% \\
    Contenu (vivier 6 h)                  & 0,0200 & \textbf{0,268 \%} & 96 \% \\
    SVD Surprise (binaire + négatifs)     & 0,0160 & 0,010 \% & 53 \% \\
    \bottomrule
  \end{tabular}

  \vspace{0.8em}
  \raggedright\small
  \begin{itemize}\itemsep3pt
    \item Retenu : \pointfort{4 populaires (1 h) + 1 contenu}. La place accordée au
          contenu ne coûte rien en précision (0,2500 contre 0,2525, écart non
          significatif) et \textbf{multiplie la couverture par 38}.
    \item Une place donnée à l'ALS coûterait 3,7 \% de précision pour un gain de
          couverture neuf fois plus faible : il reste disponible, mais n'est pas
          servi.
  \end{itemize}
\end{frame}

% ═════════════════ PARTIE 2 — Fonctionnalités (6 min, 5 diapos) ═══════════
\begin{frame}{Partie 2 — Les stratégies dans l'application}
  \framesubtitle{6 minutes \textperiodcentered\ ce que l'utilisateur peut faire}
  \small
  \begin{itemize}
    \item L'application expose les \textbf{cinq stratégies} : on bascule en direct
          et les listes diffèrent visiblement.
      \begin{itemize}
        \item « Production » — 4 populaires (1 h) + 1 contenu : ce qui est servi ;
        \item « Contenu », « ALS », « SVD Surprise », « Hybride » : pour comparaison.
      \end{itemize}
    \item Une case \pointfort{« Limiter aux articles récents »}, activée par
          défaut : la décocher montre en direct l'effet mesuré en partie 1.
    \item Le nombre d'articles est réglable de 1 à 10.
    \item Chaque article affiche sa \textbf{note en étoiles} et son nombre de
          lecteurs — les seules informations réelles disponibles sur un jeu
          anonymisé.
  \end{itemize}
\end{frame}

\begin{frame}{Un lecteur dont on ne sait rien}
  \framesubtitle{Le cas le plus fréquent sur un portail d'actualité}
  \centering
  \begin{tikzpicture}[node distance=7mm, every node/.style={font=\scriptsize},
      bloc/.style={draw=primary, thick, rounded corners, text width=2.7cm,
                   align=center, minimum height=1.1cm, text=white}]
    \node[bloc, fill=gris]    (a) {Aucun historique};
    \node[bloc, fill=accent,  right=of a] (b) {Articles récents\\les plus lus};
    \node[bloc, fill=accent2, right=of b] (c) {Région connue :\\articles récents\\de sa région};
    \draw[-{Latex[length=2mm]}, gris, thick] (a) -- (b);
    \draw[-{Latex[length=2mm]}, gris, thick] (b) -- (c);
  \end{tikzpicture}

  \vspace{1em}
  \raggedright\small
  \begin{itemize}\itemsep4pt
    \item \textbf{Cascade de repli} : région croisée avec la fraîcheur, puis
          fraîcheur, puis région, puis historique complet. Aucune requête ne
          renvoie de liste vide.
    \item La région est la seule information exploitable sur un lecteur inconnu :
          27 régions retenues, celles ayant assez de clics pour un classement.
    \item Dès le premier article lu, le profil de contenu prend le relais —
          \pointfort{sans aucun ré-entraînement}.
  \end{itemize}
\end{frame}

\begin{frame}{Inscrire un lecteur, voir son profil se former}
  \framesubtitle{En direct, sans ré-entraînement}
  \small
  \begin{itemize}
    \item L'application permet d'\textbf{inscrire un lecteur}, de lui donner un
          profil de départ, puis de marquer des articles comme lus.
    \item À chaque lecture enregistrée les recommandations changent : le profil est
          la moyenne des vecteurs lus, recalculée à la demande.
    \item Un onglet \textbf{catalogue} (364 047 articles, quatre tris, filtre par
          note) permet de cliquer autre chose que les cinq suggestions —
      \begin{itemize}
        \item sans lui, le lecteur ne pourrait cliquer que ce que le modèle propose
              déjà, et son profil ne ferait que se renforcer dans sa direction
              initiale.
      \end{itemize}
    \item Côté service, l'historique est \pointfort{transmis dans la requête} : le
          service recommande un lecteur qu'il ne connaît pas, sans conserver d'état.
  \end{itemize}
\end{frame}

\begin{frame}{Publier un nouvel article}
  \framesubtitle{Disponible en minutes, sans ré-entraînement}
  \centering
  \begin{tikzpicture}[node distance=7mm, every node/.style={font=\scriptsize},
      bloc/.style={draw=primary, thick, rounded corners, text width=2.7cm,
                   align=center, minimum height=1cm, text=white}]
    \node[bloc, fill=accent]  (a) {Embedding\\250 dimensions};
    \node[bloc, fill=accent2, right=of a] (b) {Projection ACP\\sauvegardée};
    \node[bloc, fill=primary, right=of b] (c) {Ajout au catalogue\\(50 dimensions)};
    \draw[-{Latex[length=2mm]}, gris, thick] (a) -- (b);
    \draw[-{Latex[length=2mm]}, gris, thick] (b) -- (c);
  \end{tikzpicture}

  \vspace{1em}
  \raggedright\small
  \begin{itemize}\itemsep4pt
    \item La projection ACP est sérialisée avec le catalogue (\textbf{51 Ko}) :
          sans elle, intégrer un article imposerait de réajuster l'ACP, donc de
          recalculer les 364 047 vecteurs existants.
    \item Vérifié : trois articles ajoutés reçoivent les identifiants
          364047--364049 et \textbf{entrent immédiatement} dans le top-5 d'un
          lecteur proche.
    \item L'article n'a ni étoiles ni facteurs collaboratifs — il n'a aucun clic.
          Il est recommandé par le contenu, seul signal disponible à cet instant.
  \end{itemize}
\end{frame}

\begin{frame}{Reproductibilité et garde-fous}
  \framesubtitle{Ce qui empêche une régression de partir en production}
  \centering
  \scriptsize
  \begin{tabular}{@{}ll@{}}
    \toprule
    \textbf{Élément} & \textbf{Rôle} \\
    \midrule
    Découpage temporel 60/20/20 & réglage sur validation, mesure sur test \\
    MLflow (SQLite)             & paramètres, métriques et commit git de chaque essai \\
    Registre de modèles         & versions des artefacts, retour arrière possible \\
    \textbf{Porte de qualité}   & bloque la publication si HitRate@5 baisse de $>$10 \% \\
    CI (GitHub Actions)         & 48 tests, cohérence des copies déployées \\
    \bottomrule
  \end{tabular}

  \vspace{0.8em}
  \raggedright\small
  \begin{itemize}\itemsep3pt
    \item La porte compare aux métriques de référence \textbf{versionnées dans le
          dépôt} ; la référence ne change que délibérément.
    \item Sans elle, un ré-entraînement dégradé se publierait sans alerte — c'est
          exactement ce qui est arrivé avec un chargeur d'artefacts incomplet, qui
          servait la popularité de tout l'historique \pointfort{sans lever
          d'erreur}.
  \end{itemize}
\end{frame}

% ═══════════════ PARTIE 3 — Architecture cible (2 min, 2 diapos) ══════════
\begin{frame}{Partie 3 — Architecture technique}
  \framesubtitle{2 minutes \textperiodcentered\ un cœur, trois déploiements indépendants}
  \centering
  \begin{tikzpicture}[every node/.style={font=\scriptsize, align=center, text=white},
      bloc/.style={draw=primary, thick, rounded corners, text width=3.2cm,
                   minimum height=1.3cm}]
    \node[bloc, fill=primary, text width=6cm, minimum height=0.9cm] (art)
      {Artefacts produits hors-ligne (ACP, ALS, SVD, fenêtres)};
    \node[bloc, fill=accent,  below left=1cm and 1.2cm of art] (az)
      {Azure Function\\+ Blob Storage\\\vspace{2pt}en ligne, 0,2 s};
    \node[bloc, fill=accent2, below=1cm of art] (hf)
      {Hugging Face Space\\+ HF Hub\\\vspace{2pt}démo publique};
    \node[bloc, fill=vert, below right=1cm and 1.2cm of art] (loc)
      {Application locale\\artefacts sur disque\\\vspace{2pt}aucun cloud};
    \draw[-{Latex[length=2mm]}, gris, thick] (art) -- (az);
    \draw[-{Latex[length=2mm]}, gris, thick] (art) -- (hf);
    \draw[-{Latex[length=2mm]}, gris, thick] (art) -- (loc);
  \end{tikzpicture}

  \vspace{0.6em}
  \raggedright\scriptsize
  \begin{itemize}\itemsep3pt
    \item Même cœur de recommandation dans les trois : les copies déployées sont
          \textbf{générées}, et la CI vérifie qu'elles sont à jour.
    \item À l'inférence, aucune bibliothèque d'apprentissage : \pointfort{numpy
          seul}. scikit-learn, implicit et Surprise ne servent qu'hors-ligne.
    \item Architecture 2 retenue pour le MVP : la Function lit les artefacts dans
          Blob Storage, sans API intermédiaire.
  \end{itemize}
\end{frame}

\begin{frame}{Architecture cible}
  \framesubtitle{Absorber les nouveaux articles et les nouveaux lecteurs}
  \centering
  \scriptsize
  \begin{tabular}{@{}lllc@{}}
    \toprule
    \textbf{Événement} & \textbf{Traitement} & \textbf{Latence} &
    \textbf{Ré-entr.} \\
    \midrule
    Nouvel article      & embedding + projection ACP        & minutes    & non \\
    Nouveau clic        & profil recalculé à la demande     & temps réel & non \\
    Nouveau lecteur     & popularité, puis contenu au 1	extsuperscript{er} clic
                        & temps réel & non \\
    Couverture collab.  & facteurs latents                  & heures     & oui \\
    Dérive du catalogue & ré-ajustement de l'ACP            & planifié   & oui \\
    \bottomrule
  \end{tabular}

  \vspace{0.8em}
  \raggedright\small
  \begin{itemize}\itemsep3pt
    \item \pointfort{Seule la partie collaborative impose un ré-entraînement} ; les
          deux besoins immédiats sont couverts sans.
    \item Reste à construire : ingestion événementielle, store de profils, fenêtre
          glissante recalculée en continu, index de similarité pour l'échelle.
  \end{itemize}
\end{frame}

% ═════════════════════ PARTIE 4 — Démonstration (2 min, 1 diapo) ══════════
\begin{frame}{Partie 4 — Démonstration}
  \framesubtitle{2 minutes \textperiodcentered\ six gestes}
  \begin{enumerate}
    \item Sélectionner un lecteur : 5 articles, avec étoiles et nombre de lecteurs.
    \item Décocher « articles récents » : la liste change — \pointfort{l'effet
          mesuré, en direct}.
    \item Basculer sur « SVD Surprise » : autre classement, autre bibliothèque.
    \item Parcourir le catalogue : cliquer un article hors des suggestions.
    \item Inscrire un lecteur : il reçoit des recommandations immédiatement.
    \item Le même appel sur Azure : 0,2 s, réponse identique au local.
  \end{enumerate}
  \aretenir{Les articles n'ont ni titre ni texte : le jeu de données est anonymisé
    par sa source. L'interface affiche identifiant, catégorie, longueur, date et
    note — les seules informations réelles disponibles.}
\end{frame}

% ═══════════════════════════════════════════════════════════ Clôture ═══════
{\setbeamercolor{background canvas}{bg=primary}
\begin{frame}[plain]
  \vfill
  {\color{white}\LARGE\bfseries La fraîcheur pèse plus que le modèle}\\[1em]
  {\color{white!75!primary}\small
   Facteur 250 sur la précision, mesuré sur découpage temporel.\\[0.4em]
   Une place sur cinq accordée à la personnalisation ne coûte rien
   et multiplie par 38 la part du catalogue exposée.}
  \vfill
\end{frame}}

\end{document}
